\documentclass[12pt]{article}

% ============================================================
%  PAQUETES
% ============================================================
\usepackage[spanish]{babel}

% ------------------------------------------------------------
%  TIPOGRAFÍA: Arial en el texto, monoespaciada en el código
%  IMPORTANTE: este archivo se compila con XeLaTeX o LuaLaTeX,
%  no con pdfLaTeX. En Overleaf: Menu > Compiler > XeLaTeX.
%  Si Arial no está instalado se usa Liberation Sans, que tiene
%  exactamente las mismas métricas y no altera la paginación.
% ------------------------------------------------------------
\usepackage{fontspec}
\IfFontExistsTF{Arial}{%
  \setmainfont{Arial}%
  \setsansfont{Arial}%
}{%
  \setmainfont{Liberation Sans}%
  \setsansfont{Liberation Sans}%
}
\IfFontExistsTF{Consolas}{%
  \setmonofont{Consolas}[Scale=MatchLowercase]%
}{%
  \setmonofont{Liberation Mono}[Scale=MatchLowercase]%
}

\usepackage[a4paper, margin=2.5cm]{geometry}
\usepackage{graphicx}          % Imágenes
\usepackage{float}             % Anclaje de figuras con [H]
\usepackage{enumitem}          % Control de listas
\usepackage{needspace}         % Evita separar un bloque de su captura
\usepackage{tocloft}           % Ancho de la numeración en el índice
\setlength{\cftsecnumwidth}{2em}
\setlength{\cftsubsecnumwidth}{2.8em}
\usepackage{xcolor}            % Colores
\usepackage{listings}          % Bloques de código
\usepackage{fancyhdr}          % Encabezado y pie de página
\usepackage{titlesec}          % Estilo de títulos de sección
\usepackage{microtype}         % Mejor justificación / tipografía
\usepackage{setspace}          % Interlineado
\onehalfspacing                % Interlineado 1.5 en todo el documento
\usepackage[font=small, labelfont={bf,color=purple-uach}]{caption}
\usepackage[hidelinks]{hyperref} % Enlaces e índice clicable
\usepackage{xurl}
\urlstyle{same}

\captionsetup{font={small,stretch=1.0}}    % pies de figura a renglón sencillo

\graphicspath{{img/}{./}}

% ============================================================
%  COLORES
% ============================================================
\definecolor{purple-uach}{HTML}{371e49}

% ============================================================
%  DATOS DE LA TAREA
% ============================================================
\newcommand{\alumnoUno}{Osvaldo Hernández Juárez}
\newcommand{\matriculaUno}{385493}
\newcommand{\alumnoDos}{Diego Flores Verdad Grijalva}
\newcommand{\matriculaDos}{385660}
\newcommand{\alumnoTres}{Luis Felipe Domínguez Chávez}
\newcommand{\matriculaTres}{385500}
\newcommand{\materia}{Bases de Datos Avanzadas}
\newcommand{\docente}{José Saúl De Lira Miramontes}
\newcommand{\grupoCarrera}{Grupo 5k2}
\newcommand{\actividad}{Proyecto Final: Aplicación de Base de Datos de Grafos}
% Versión de la portada, con el salto de renglón donde se ve mejor.
\newcommand{\actividadPortada}{Proyecto Final:\\ Blog en Neo4j}
\newcommand{\fechaEntrega}{24 de noviembre de 2026}

\hypersetup{
    pdftitle={\actividad},
    pdfauthor={\alumnoUno, \alumnoDos, \alumnoTres},
    pdfsubject={\materia}
}

% ============================================================
%  ESTILO DE TÍTULOS DE SECCIÓN
% ============================================================
\titleformat{\section}
  {\large\bfseries\color{purple-uach}}{\thesection}{1em}{}
\titleformat{\subsection}
  {\normalsize\bfseries\color{purple-uach}}{\thesubsection}{1em}{}
\titlespacing*{\section}{0pt}{2.4ex plus 1ex minus .2ex}{1.2ex plus .2ex}

% ============================================================
%  PIE DE PÁGINA
% ============================================================
\pagestyle{fancy}
\fancyhf{}
\renewcommand{\headrulewidth}{0pt}
\renewcommand{\footrulewidth}{0pt}
\fancyfoot[R]{\thepage}

% ============================================================
%  ESTILO DE BLOQUES DE CÓDIGO
% ============================================================
\lstdefinelanguage{mongosh}{
    morekeywords={db,sh,rs,const,let,new,find,findOne,insertOne,insertMany,updateOne,updateMany,deleteOne,deleteMany,countDocuments,aggregate,runCommand,getSiblingDB,createCollection,drop,print,printjson},
    morestring=[b]",
    morestring=[b]',
    morecomment=[l]{//},
    sensitive=true
}

\lstdefinestyle{codeStyle}{
    backgroundcolor=\color{gray!8},
    commentstyle=\color{purple-uach!70},
    keywordstyle=\color{purple-uach}\bfseries,
    stringstyle=\color{teal},
    basicstyle=\setstretch{1}\ttfamily\footnotesize,
    breaklines=true,
    breakatwhitespace=false,
    frame=single,
    rulecolor=\color{gray!40},
    tabsize=2,
    captionpos=b,
    aboveskip=1.2ex,
    belowskip=1.2ex
}
\lstset{style=codeStyle, language=mongosh}
\lstdefinestyle{pyStyle}{style=codeStyle, language=Python, showstringspaces=false, morekeywords={self,None,True,False}}
\lstdefinestyle{plainStyle}{style=codeStyle, language={}, showstringspaces=false}


% ============================================================
%  CAPTURAS Y PENDIENTES
% ============================================================
% \captura{archivo}{pie de figura}{ancho}: si la imagen todavía no está en
% img/, se muestra un recuadro con el nombre del archivo que falta.
\newcommand{\captura}[3]{%
  \begin{figure}[H]
    \centering
    \IfFileExists{img/#1}{%
      \includegraphics[width=#3\linewidth]{#1}%
    }{%
      \fbox{\parbox[c][2.5cm][c]{0.8\linewidth}{\centering\small\color{gray}Falta la captura \texttt{img/#1}}}%
    }
    \caption{#2}
  \end{figure}}
% \pendiente{texto}: aviso visible de algo que falta; quitarlo antes de entregar.
\newcommand{\pendiente}[1]{%
  \par\noindent\fcolorbox{purple-uach}{purple-uach!6}{\parbox{\dimexpr\linewidth-2\fboxsep-2\fboxrule}{\small\textbf{Pendiente:} #1}}\par}

\lstdefinelanguage{cypher}{
    morekeywords={MATCH,OPTIONAL,WHERE,RETURN,CREATE,MERGE,SET,DELETE,DETACH,WITH,UNWIND,AS,ORDER,BY,DESC,ASC,DISTINCT,AND,OR,NOT,IN,EXISTS,COUNT,CONSTRAINT,FOR,REQUIRE,IS,UNIQUE,IF,CALL,YIELD,LIMIT,SHOW},
    sensitive=true, morecomment=[l]{//}, morestring=[b]", morestring=[b]'}
\lstdefinestyle{cypherStyle}{style=codeStyle, language=cypher, showstringspaces=false}
\lstdefinestyle{jsStyle}{style=codeStyle, language=mongosh, showstringspaces=false, morecomment=[s]{/*}{*/}, morestring=[b]`}

\emergencystretch=1.5em

% ============================================================
\begin{document}

% ---------------- PORTADA ----------------
\begin{titlepage}
\singlespacing
\centering
\noindent
\begin{minipage}[c]{0.5\linewidth}
    \raggedright
    \includegraphics[height=4.5cm]{logo-uach.png}
\end{minipage}%
\begin{minipage}[c]{0.5\linewidth}
    \raggedleft
    \includegraphics[height=4.5cm]{logo-fing.png}
\end{minipage}\\[1.2cm]
{\Large \textbf{Universidad Autónoma de Chihuahua}}\\[0.2cm]
{\large Facultad de Ingeniería}\\[0.2cm]
{\normalsize \grupoCarrera}\\[2cm]

{\color{purple-uach}\rule{\linewidth}{0.6mm}}\\[0.5cm]
{\huge \bfseries \linespread{1.25}\selectfont \actividadPortada\par}\vspace{0.5cm}
{\color{purple-uach}\rule{\linewidth}{0.6mm}}\\[1.5cm]

{\Large \materia}\\[0.3cm]
{\large Docente: \docente}

\vfill

\begin{flushright}
    \large
    \textbf{\alumnoUno}\\
    Matrícula: \matriculaUno\\[0.4cm]
    \textbf{\alumnoDos}\\
    Matrícula: \matriculaDos\\[0.4cm]
    \textbf{\alumnoTres}\\
    Matrícula: \matriculaTres
\end{flushright}

\vspace{1cm}
{\large \fechaEntrega}

\end{titlepage}


% Codigo de un archivo del repositorio; si todavia no existe, se avisa.
\newcommand{\codigo}[2]{%
  \IfFileExists{#2}{\lstinputlisting[style=#1]{#2}}{\pendiente{falta el archivo \texttt{\detokenize{#2}}.}}}

% ---------------- ÍNDICE ----------------
\tableofcontents
\newpage

% ============================================================
\section{Introducción}

En las actividades anteriores se trabajó con MongoDB, una base de datos de documentos. En este proyecto se usó una base de datos de grafos, Neo4j, en la que los datos se guardan como nodos y las relaciones entre ellos se almacenan explícitamente. Se desarrolló una aplicación en Python con interfaz gráfica que administra un blog: usuarios, artículos, comentarios, tags y categorías, con las relaciones del modelo indicado en la actividad. La aplicación permite crear, consultar, modificar y eliminar cada tipo de nodo, y realizar consultas que recorren el grafo.

El código de la aplicación y los archivos Cypher se encuentran en el repositorio \url{https://github.com/LF-Dguez/blog_neo4j}.

% ============================================================
\section{Objetivo}

Desarrollar una aplicación de base de datos de grafos.

% ============================================================
\section{Entorno de trabajo}

El proyecto se realizó en Windows 11 con Neo4j Desktop 2 y una instancia local de Neo4j 2026.10, a la que la aplicación se conecta por el protocolo Bolt en \texttt{localhost:7687}. La aplicación se escribió en Python 3.14 con el driver oficial de Neo4j (paquete \texttt{neo4j}) y la interfaz con Tkinter, que viene incluido en Python.

% ============================================================
\section{Modelo de datos del blog}

El modelo de la actividad tiene cinco tipos de entidad. En Neo4j cada una es una etiqueta de nodo, y cada relación del diagrama es un tipo de relación con dirección:

\begin{lstlisting}[style=plainStyle]
(:User {id, name, email})
(:Article {id, title, date, text})
(:Comment {id, name, url, text, date})
(:Tag {id, name, url})
(:Category {id, name, url})

(:User)-[:WROTE]->(:Article)            1 a N
(:User)-[:POSTED]->(:Comment)           1 a N
(:Article)-[:HAS_COMMENT]->(:Comment)   1 a N
(:Article)-[:TAGGED_WITH]->(:Tag)       N a N
(:Article)-[:IN_CATEGORY]->(:Category)  N a N
\end{lstlisting}

Las decisiones principales del modelo fueron las siguientes:

\begin{itemize}[itemsep=2pt, topsep=4pt]
    \item Cada nodo tiene una propiedad \texttt{id} generada con \texttt{randomUUID()}. La aplicación busca, modifica y elimina siempre por ese identificador.
    \item Las relaciones N a N no necesitan una tabla intermedia como en una base relacional: un artículo se conecta directamente con cada uno de sus tags y categorías.
    \item En el diagrama, los comentarios tienen las propiedades name y url. Se interpretaron como el título corto del comentario y su enlace permanente, y se agregaron text, con el contenido, y date. La url la genera la aplicación a partir del título del artículo.
    \item Las fechas de artículos y comentarios se guardan con el tipo \texttt{date} de Neo4j.
    \item Al eliminar, un usuario que tiene artículos o comentarios no se puede borrar; un artículo se elimina junto con sus comentarios; un tag, una categoría o un comentario se eliminan junto con sus relaciones.
\end{itemize}

% ============================================================
\section{Estructura del proyecto}

\begin{lstlisting}[style=plainStyle]
blog_neo4j/
  app/main.py            ventana principal con una pestaña por tipo de nodo
  app/conexion.py        conexión con Neo4j (lee config.ini)
  app/ui_comun.py        componentes de interfaz comunes
  app/cargar_cypher.py   ejecuta archivos .cypher
  app/usuarios.py        pestaña Usuarios
  app/articulos.py       pestaña Artículos
  app/comentarios.py     pestaña Comentarios
  app/tags.py            pestaña Tags
  app/categorias.py      pestaña Categorías
  app/consultas.py       pestaña Consultas
  cypher/01_restricciones.cypher ... 06_consultas.cypher
  config.ejemplo.ini     plantilla de config.ini
\end{lstlisting}

Los datos de conexión (URI, usuario y contraseña) se leen de \texttt{config.ini}, que no se sube al repositorio; en él solo está la plantilla \texttt{config.ejemplo.ini}.

% ============================================================
\section{Restricciones y datos iniciales}

\subsection{Instancia de Neo4j}

\captura{01-neo4j-desktop.png}{Instancia de Neo4j en Neo4j Desktop.}{1.0}

\subsection{Restricciones de unicidad}

Cada \texttt{id} es único, igual que el email de un usuario y el nombre de un tag o de una categoría. Una restricción \texttt{UNIQUE} también crea un índice, así que las búsquedas por esas propiedades son rápidas.

\codigo{cypherStyle}{../cypher/01_restricciones.cypher}

\subsection{Datos iniciales}

El archivo \texttt{02\_datos.cypher} borra los nodos del blog y vuelve a crear 5 usuarios, 4 categorías, 8 tags, 8 artículos y 14 comentarios con todas sus relaciones. Se usa \texttt{UNWIND} para crear varios nodos a partir de una lista en una sola sentencia.

\codigo{cypherStyle}{../cypher/02_datos.cypher}

\subsection{Cargador de archivos Cypher}

\texttt{cargar\_cypher.py} ejecuta en orden las sentencias de uno o varios archivos \texttt{.cypher}. De cada una muestra su título (el comentario que tiene arriba), lo que cambió en la base y, si tiene \texttt{RETURN}, sus primeras filas. Al final cuenta los nodos y las relaciones del blog.

\begin{lstlisting}[style=plainStyle]
python app/cargar_cypher.py cypher/01_restricciones.cypher cypher/02_datos.cypher
\end{lstlisting}

\captura{02-cargar-datos.png}{Creación de las restricciones y carga de los datos iniciales.}{1.0}

Se crearon las 8 restricciones y los 39 nodos, con 65 relaciones en total.

\captura{03-restricciones.png}{Restricciones de la base de datos en Neo4j Desktop.}{1.0}

\captura{04-grafo-completo.png}{Grafo completo del blog.}{1.0}

% ============================================================
\section{Conexión y ventana principal}

\subsection{Conexión con Neo4j}

La clase \texttt{BaseDatos} lee \texttt{config.ini}, crea el driver de Neo4j y ofrece tres métodos: \texttt{consultar} para lecturas, \texttt{ejecutar} para escrituras y \texttt{transaccion} para varias sentencias que deben aplicarse juntas. Los valores siempre se envían como parámetros (\texttt{\$id}, \texttt{\$name}), nunca pegados al texto de la consulta, y los tipos de Neo4j como \texttt{date} se convierten a texto para mostrarlos.

\codigo{pyStyle}{../app/conexion.py}

\subsection{Componentes comunes de la interfaz}

Todas las pestañas de administración tienen la misma distribución: a la izquierda una tabla con los registros y a la derecha un formulario con los botones Nuevo, Guardar y Eliminar. Esa distribución está en la clase \texttt{PestanaCRUD}; cada pestaña hereda de ella y solo define sus columnas, sus campos y el Cypher de cada operación. Para las relaciones hay dos controles: \texttt{Selector}, una lista desplegable para elegir un nodo (por ejemplo, el autor), y \texttt{ListaMultiple}, para elegir varios (por ejemplo, los tags).

\codigo{pyStyle}{../app/ui_comun.py}

\subsection{Ventana principal}

\texttt{main.py} se conecta a la base, muestra la versión del servidor en el encabezado y agrega una pestaña por cada módulo. Al cambiar de pestaña se recargan sus datos, porque otra pestaña pudo haber creado o eliminado nodos relacionados.

\codigo{pyStyle}{../app/main.py}

% ============================================================
\section{Usuarios}

\codigo{pyStyle}{../app/usuarios.py}

La tabla muestra cuántos artículos y comentarios tiene cada usuario con \texttt{COUNT \{ \}}, que cuenta los patrones de relaciones de cada nodo sin agruparlos.

\captura{05-app-usuarios.png}{Pestaña Usuarios.}{1.0}

\captura{06-usuario-crear.png}{Creación de un usuario.}{1.0}

\captura{07-usuario-editar.png}{Modificación de un usuario.}{1.0}

\captura{08-usuario-email-repetido.png}{Intento de crear un usuario con un email repetido.}{1.0}

La restricción \texttt{user\_email} impide que dos usuarios tengan el mismo email; Neo4j rechaza la escritura y la aplicación muestra el error.

\captura{09-usuario-eliminar-bloqueado.png}{Intento de eliminar un usuario con artículos y comentarios.}{1.0}

\captura{10-usuario-eliminar.png}{Eliminación de un usuario.}{1.0}

Las mismas operaciones, escritas en Cypher con un usuario de prueba:

\codigo{cypherStyle}{../cypher/03_usuarios.cypher}

\captura{11-cypher-03.png}{Ejecución de las operaciones de usuarios en Cypher.}{1.0}

% ============================================================
\section{Artículos y comentarios}

\subsection{Artículos}

\codigo{pyStyle}{../app/articulos.py}

La pestaña Artículos administra los nodos \texttt{:Article} y, además de sus propiedades, las tres relaciones que salen o llegan a un artículo: el autor (\texttt{WROTE}), los tags (\texttt{TAGGED\_WITH}) y las categorías (\texttt{IN\_CATEGORY}). El autor se elige con un \texttt{Selector}, porque un artículo tiene un solo autor, y los tags y las categorías con dos \texttt{ListaMultiple}, porque un artículo puede tener varios.

Para llenar la tabla, \texttt{listar()} usa \textit{pattern comprehensions}, que recorren las relaciones de cada artículo y regresan una lista; por ejemplo, la de los nombres de sus tags:

\begin{lstlisting}[style=cypherStyle]
[(a)-[:TAGGED_WITH]->(t:Tag) | t.name] AS tags
\end{lstlisting}

\noindent Si se usara un \texttt{MATCH} por cada relación, cada artículo aparecería repetido una vez por cada combinación de tag y categoría. El número de comentarios se obtiene con \texttt{COUNT \{ \}}.

Crear o modificar un artículo requiere varias sentencias: la del nodo, la de su autor y las de sus tags y categorías. Todas se ejecutan en una sola transacción con \texttt{BaseDatos.transaccion()}, de modo que si una falla no se guarda ninguna. Para actualizar los tags y las categorías se borran las relaciones anteriores y se crean las que están elegidas en el formulario, lo que cubre tanto agregar como quitar. Por esa misma razón el \texttt{id} de un artículo nuevo se genera en Python con \texttt{uuid4()}: las sentencias que crean las relaciones lo necesitan dentro de la misma transacción.

Al eliminar un artículo también se eliminan sus comentarios, porque un comentario no tiene sentido sin el artículo al que pertenece; \texttt{DETACH DELETE} borra los nodos junto con todas sus relaciones.

\captura{20-articulos-lista.png}{Pestaña Artículos.}{1.0}
\captura{21-articulo-crear.png}{Creación de un artículo con autor, tags y categoría.}{1.0}
\captura{22-articulo-editar.png}{Modificación de los tags y categorías de un artículo.}{1.0}

\subsection{Comentarios}

\codigo{pyStyle}{../app/comentarios.py}

La pestaña Comentarios administra los nodos \texttt{:Comment}. Cada comentario se conecta con dos nodos: el usuario que lo publicó (\texttt{POSTED}) y el artículo al que pertenece (\texttt{HAS\_COMMENT}); ambos se eligen con un \texttt{Selector} y son obligatorios.

La propiedad \texttt{url} no se captura en el formulario: la genera Cypher al guardar, a partir del título del artículo en minúsculas y con guiones, más los primeros ocho caracteres del \texttt{id} del comentario. Así cada comentario tiene un enlace permanente único, por ejemplo \texttt{/articulos/grafos-en-la-vida-diaria\#comentario-d6c2cc67}. Al modificar un comentario se reemplazan sus dos relaciones y se vuelve a calcular la \texttt{url}, por si se cambió de artículo.

Eliminar un comentario solo borra ese nodo y sus relaciones; el usuario y el artículo se conservan. En las capturas se crea un comentario sobre el artículo de la sección anterior, se modifica, se comprueba que la columna \textit{Coment.} del artículo cambia a 1 y después se eliminan el comentario y el artículo, con lo que la base regresa a sus 8 artículos y 14 comentarios.

\captura{23-comentarios-lista.png}{Pestaña Comentarios.}{1.0}
\captura{24-comentario-crear.png}{Creación de un comentario.}{1.0}
\captura{25-comentario-editar.png}{Modificación de un comentario.}{1.0}
\captura{26-articulo-con-comentario.png}{Artículo con su nuevo comentario.}{1.0}
\captura{27-comentario-eliminar.png}{Eliminación de un comentario.}{1.0}
\captura{28-articulo-eliminar.png}{Eliminación de un artículo.}{1.0}

\subsection{Operaciones en Cypher}

Las mismas operaciones se escribieron en Cypher con un artículo y un comentario de prueba. El archivo se ejecuta en orden: crea, lee y modifica el artículo, hace lo mismo con el comentario, los elimina y al final comprueba que no quedó ningún nodo de prueba.

\codigo{cypherStyle}{../cypher/04_articulos_comentarios.cypher}

\captura{29-cypher-04.png}{Ejecución de las operaciones de artículos y comentarios en Cypher.}{1.0}
\captura{30-grafo-articulo.png}{Un artículo con su autor, tags, categorías y comentarios.}{1.0}

% ============================================================
\section{Tags y categorías}

\subsection{Tags}

\codigo{pyStyle}{../app/tags.py}

Tags y categorías son nodos sencillos, con nombre y url, que se conectan con los artículos mediante relaciones N a N. Por eso sus pestañas solo definen columnas, campos y el Cypher de cada operación, y el resto lo hereda de \texttt{PestanaCRUD}. La columna \textit{Artículos} cuenta con \texttt{COUNT \{ \}} cuántos artículos llegan a cada tag.

El nombre de un tag no se puede repetir gracias a la restricción \texttt{tag\_name}: al intentar crear otro tag \texttt{neo4j}, Neo4j rechaza la escritura y la aplicación muestra el error. Al eliminar un tag se usa \texttt{DETACH DELETE}, que borra el nodo y sus relaciones \texttt{TAGGED\_WITH}, pero no los artículos.

\captura{40-tags-lista.png}{Pestaña Tags.}{1.0}
\captura{41-tag-crear.png}{Creación de un tag.}{1.0}
\captura{42-tag-editar.png}{Modificación de un tag.}{1.0}
\captura{43-tag-repetido.png}{Intento de crear un tag repetido.}{1.0}
\captura{44-tag-eliminar.png}{Eliminación de un tag.}{1.0}

\subsection{Categorías}

\codigo{pyStyle}{../app/categorias.py}

La pestaña Categorías funciona igual que la de Tags, con la restricción \texttt{category\_\allowbreak name} para el nombre y la relación \texttt{IN\_CATEGORY}. Al eliminar una categoría, los artículos que la tenían se conservan con sus demás categorías.

\captura{45-categorias-lista.png}{Pestaña Categorías.}{1.0}
\captura{46-categoria-crear.png}{Creación de una categoría.}{1.0}
\captura{47-categoria-editar.png}{Modificación de una categoría.}{1.0}
\captura{48-categoria-eliminar.png}{Eliminación de una categoría.}{1.0}

\subsection{Operaciones en Cypher}

Las operaciones de tags y categorías en Cypher usan un tag y una categoría de prueba, que se crean, se consultan junto con el número de artículos de cada uno, se modifican y se eliminan; la última sentencia comprueba que no quedó ningún nodo de prueba.

\codigo{cypherStyle}{../cypher/05_tags_categorias.cypher}

\captura{57-cypher-05.png}{Ejecución de las operaciones de tags y categorías en Cypher.}{1.0}

% ============================================================
\section{Consultas sobre el grafo}

\codigo{pyStyle}{../app/consultas.py}

La pestaña Consultas no modifica la base: reúne ocho consultas de solo lectura que recorren el grafo. Se elige una de la lista y, si necesita un parámetro (un tag o un artículo), se elige en un \texttt{Selector}; el resultado se muestra en una tabla cuyas columnas son las del \texttt{RETURN}. Estas consultas son las que muestran la ventaja de un grafo, porque siguen varias relaciones seguidas sin necesidad de \texttt{JOIN}.

\begin{enumerate}[itemsep=3pt, topsep=4pt]
    \item \textbf{Artículos por categoría.} Usa \texttt{OPTIONAL MATCH} para que aparezcan también las categorías sin artículos. Bases de Datos tiene 6 artículos, Programación 3, y Ciencia de Datos e Infraestructura 1 cada una.
    \item \textbf{Tags más usados.} Cuenta con \texttt{COUNT \{ \}} los artículos de cada tag: neo4j aparece en 4, cypher y grafos en 3, mongodb, nosql y python en 2, y docker y sql en 1.
    \item \textbf{Artículos más comentados.} Cada artículo con su autor y su número de comentarios. El primero es ``Introducción a las bases de datos de grafos'', con 3, y el último ``De SQL a Cypher: guía rápida'', con 0.
    \item \textbf{Actividad de cada usuario.} Artículos escritos y comentarios publicados por cada usuario, ordenados por la suma de ambos; por ejemplo, Ana Torres tiene 2 artículos y 3 comentarios, y Carlos Méndez 1 artículo y 4 comentarios.
    \item \textbf{Artículos de un tag.} Recibe el tag como parámetro. Con \texttt{neo4j} regresa 4 artículos, del más reciente (``Conectar Python con Neo4j'') al más antiguo.
    \item \textbf{Artículos relacionados.} A partir de un artículo recorre el patrón artículo--tag--artículo para encontrar los que comparten tags, y cuenta cuántos tienen en común. Para ``Introducción a las bases de datos de grafos'', el más relacionado es ``MongoDB vs Neo4j: cuándo usar cada uno'', con 2 tags en común, y otros cuatro artículos comparten uno.
    \item \textbf{Quién comenta a quién.} Recorre usuario--comentario--artículo--autor y descarta a quien comenta su propio artículo con \texttt{WHERE lector <> autor}. Regresa 12 pares; los más frecuentes son Carlos Méndez comentando a Ana Torres y a Sofía Chávez, 2 veces cada uno.
    \item \textbf{Artículos sin comentarios.} Usa \texttt{WHERE NOT EXISTS \{ \}} para quedarse con los artículos que no tienen ninguna relación \texttt{HAS\_COMMENT}. Solo ``De SQL a Cypher: guía rápida'', de Luis Ramírez, cumple la condición.
\end{enumerate}

\captura{49-consulta-1.png}{Artículos por categoría.}{1.0}
\captura{50-consulta-2.png}{Tags más usados.}{1.0}
\captura{51-consulta-3.png}{Artículos más comentados.}{1.0}
\captura{52-consulta-4.png}{Actividad de cada usuario.}{1.0}
\captura{53-consulta-5.png}{Artículos de un tag.}{1.0}
\captura{54-consulta-6.png}{Artículos relacionados por tags.}{1.0}
\captura{55-consulta-7.png}{Quién comenta a quién.}{1.0}
\captura{56-consulta-8.png}{Artículos sin comentarios.}{1.0}

Las mismas consultas están en \texttt{06\_consultas.cypher} para ejecutarlas fuera de la aplicación; en lugar del parámetro usan valores de ejemplo (el tag \texttt{neo4j} y el artículo ``Introducción a las bases de datos de grafos'').

\codigo{cypherStyle}{../cypher/06_consultas.cypher}

\captura{58a-cypher-06.png}{Ejecución de las consultas en Cypher (primera parte).}{1.0}
\captura{58b-cypher-06.png}{Ejecución de las consultas en Cypher (segunda parte).}{1.0}
\captura{59-grafo-tags.png}{Artículos y sus tags.}{1.0}

% ============================================================
\section{Conclusión}

Se desarrolló una aplicación en Python que administra un blog guardado en Neo4j, con el modelo de usuarios, artículos, comentarios, tags y categorías indicado en la actividad. Cada tipo de nodo tiene su pestaña con altas, consultas, modificaciones y bajas, incluidas sus relaciones, y una pestaña adicional reúne consultas que recorren el grafo.

El proyecto mostró las diferencias de trabajar con una base de datos de grafos frente a las bases de documentos de las actividades anteriores. Las relaciones son datos de primera clase: una relación N a N, como la de artículos y tags, no necesita una tabla o colección intermedia, y consultas como ``artículos relacionados'' o ``quién comenta a quién'' se escriben como un patrón que sigue varias relaciones, en lugar de varios \texttt{JOIN} o \textit{lookups}. Las restricciones de unicidad mantuvieron la integridad de los datos, las transacciones permitieron guardar un nodo junto con todas sus relaciones de forma atómica, y el uso de parámetros en todo el Cypher evitó armar consultas pegando texto.

Las capturas muestran que cada operación deja la base en el estado esperado: los conteos de nodos y relaciones regresan a sus valores iniciales después de crear y eliminar los registros de prueba, tanto desde la aplicación como desde los archivos Cypher.

% ============================================================
\section*{Referencias}
\addcontentsline{toc}{section}{Referencias}

\begin{list}{}{%
    \setlength{\leftmargin}{1.27cm}
    \setlength{\itemindent}{-1.27cm}
    \setlength{\itemsep}{6pt}
    \setlength{\parsep}{0pt}}
\item Neo4j, Inc. (s.f.). \textit{Cypher Manual}. \url{https://neo4j.com/docs/cypher-manual/current/}
\item Neo4j, Inc. (s.f.). \textit{Neo4j Python Driver Manual}. \url{https://neo4j.com/docs/python-manual/current/}
\item Python Software Foundation. (s.f.). \textit{tkinter — Python interface to Tcl/Tk}. \url{https://docs.python.org/3/library/tkinter.html}
\end{list}

\end{document}
